\documentclass[tikz,border=10pt]{standalone}

\usepackage{amsmath,amssymb}
\usepackage{fontspec}
\usepackage{xeCJK}
\IfFontExistsTF{SimHei}{\setCJKmainfont{SimHei}}{\setCJKmainfont{Noto Sans CJK SC}}
\usepackage{xcolor}
\usepackage{tikz}
\usetikzlibrary{arrows.meta,backgrounds,fit,positioning}

\definecolor{ink}{HTML}{18324A}
\definecolor{muted}{HTML}{5F7286}
\definecolor{linec}{HTML}{C9D8E6}
\definecolor{panel}{HTML}{F8FBFD}
\definecolor{bluefill}{HTML}{EFF6FF}
\definecolor{blueedge}{HTML}{9FC2F1}
\definecolor{greenfill}{HTML}{EFFAF4}
\definecolor{greenedge}{HTML}{A8D7B6}
\definecolor{purplefill}{HTML}{F5F1FC}
\definecolor{purpleedge}{HTML}{CBB9E7}
\definecolor{orangefill}{HTML}{FFF6ED}
\definecolor{orangeedge}{HTML}{F2C892}
\definecolor{tealfill}{HTML}{EEF9F8}
\definecolor{tealedge}{HTML}{9FD9CF}
\definecolor{badgeblue}{HTML}{2F6BFF}
\definecolor{badgegreen}{HTML}{1B9E77}
\definecolor{badgeorange}{HTML}{E17A21}
\definecolor{badgepurple}{HTML}{6A4C93}

\tikzset{
  stepbox/.style={
    rectangle,
    rounded corners=9pt,
    draw=linec,
    fill=panel,
    minimum width=13.7cm,
    minimum height=1.85cm,
    text width=13.1cm,
    inner sep=7pt,
    align=left,
    font=\sffamily\small
  },
  stagepill/.style={
    rectangle,
    rounded corners=9pt,
    fill=#1,
    inner sep=5pt,
    text=white,
    font=\sffamily\scriptsize\bfseries
  },
  stagepill/.default=badgeblue,
  idbubble/.style={
    circle,
    fill=#1,
    minimum size=0.95cm,
    text=white,
    font=\sffamily\small\bfseries,
    inner sep=0pt
  },
  idbubble/.default=badgeblue,
  steptext/.style={font=\sffamily\small\bfseries, text=ink, align=left},
  note/.style={font=\sffamily\scriptsize, text=muted, align=left},
  resultchip/.style={
    rectangle,
    rounded corners=8pt,
    fill=#1,
    inner sep=5pt,
    text=white,
    font=\sffamily\tiny\bfseries,
    align=center
  },
  resultchip/.default=badgegreen,
  link/.style={-{Stealth[length=3.4mm]}, line width=0.95pt, draw=muted},
}

\begin{document}
\begin{tikzpicture}[node distance=0.8cm]

\node[stagepill] (title) at (0,0) {TVM 主线证据链};
\node[note, text width=13.3cm, below=0.55cm of title] (subtitle)
{图中只保留真正决定主线结果的五个阶段：先解决“能不能跑”，再收敛目标与调优预算，最后回到真实重建与热点算子整合。};

\node[idbubble=badgeblue, below=1.0cm of subtitle, xshift=-6.9cm] (n1) {1};
\node[stepbox, right=0.55cm of n1] (s1) {};
\node[stagepill=badgeblue, anchor=north west] at ([xshift=0.28cm,yshift=-0.28cm]s1.north west) {安全运行时重建};
\node[steptext, anchor=north west] at ([xshift=0.32cm,yshift=-0.92cm]s1.north west)
{剥离旧依赖并恢复 TVM 0.24 safe runtime，解除飞腾端加载阶段的 SIGILL / 兼容性阻塞。};
\node[resultchip=badgeblue, anchor=north east] at ([xshift=-0.32cm,yshift=-0.32cm]s1.north east) {运行时可稳定加载};

\node[idbubble=badgepurple, below=1.08cm of n1] (n2) {2};
\node[stepbox, right=0.55cm of n2] (s2) {};
\node[stagepill=badgepurple, anchor=north west] at ([xshift=0.28cm,yshift=-0.28cm]s2.north west) {目标收敛};
\node[steptext, anchor=north west] at ([xshift=0.32cm,yshift=-0.92cm]s2.north west)
{把 target 从 generic 路线收敛到 \texttt{cortex-a72 + neon}，确保 artifact 与飞腾微架构一致。};
\node[resultchip=badgepurple, anchor=north east] at ([xshift=-0.32cm,yshift=-0.32cm]s2.north east) {generic $\rightarrow$ A72};

\node[idbubble=badgegreen, below=1.08cm of n2] (n3) {3};
\node[stepbox, right=0.55cm of n3] (s3) {};
\node[stagepill=badgegreen, anchor=north west] at ([xshift=0.28cm,yshift=-0.28cm]s3.north west) {增量调优};
\node[steptext, anchor=north west] at ([xshift=0.32cm,yshift=-0.92cm]s3.north west)
{先用 rebuild-only / payload-only 缩小搜索面，再把 500 次 trial 集中投入真实瓶颈。};
\node[resultchip=badgegreen, anchor=north east] at ([xshift=-0.32cm,yshift=-0.32cm]s3.north east) {payload 152.8 ms};

\node[idbubble=badgeorange, below=1.08cm of n3] (n4) {4};
\node[stepbox, right=0.55cm of n4] (s4) {};
\node[stagepill=badgeorange, anchor=north west] at ([xshift=0.28cm,yshift=-0.28cm]s4.north west) {真实重建评测};
\node[steptext, anchor=north west] at ([xshift=0.32cm,yshift=-0.92cm]s4.north west)
{回到 300 张图像的 latent$\rightarrow$PNG 端到端链路，确认收益不是只存在于 micro benchmark。};
\node[resultchip=badgeorange, anchor=north east] at ([xshift=-0.32cm,yshift=-0.32cm]s4.north east) {端到端 230.3 ms};

\node[idbubble=badgeblue, below=1.08cm of n4] (n5) {5};
\node[stepbox, right=0.55cm of n5] (s5) {};
\node[stagepill=badgeblue, anchor=north west] at ([xshift=0.28cm,yshift=-0.28cm]s5.north west) {Profiling 与手写算子整合};
\node[steptext, anchor=north west] at ([xshift=0.32cm,yshift=-0.92cm]s5.north west)
{按热点证据链保留正收益 lane，并与 big.LITTLE 流水线结果一起回收为系统主线。};
\node[resultchip=badgegreen, anchor=north east] at ([xshift=-0.32cm,yshift=-0.32cm]s5.north east) {流水线 134.6 ms};

\draw[link] (n1.south) -- (n2.north);
\draw[link] (n2.south) -- (n3.north);
\draw[link] (n3.south) -- (n4.north);
\draw[link] (n4.south) -- (n5.north);

\node[
  rectangle,
  rounded corners=9pt,
  draw=linec,
  fill=bluefill,
  text width=14.7cm,
  inner sep=8pt,
  align=left,
  font=\sffamily\scriptsize,
  text=muted,
  below=0.9cm of s5
] (footer)
{读图方式：这张图只讲 TVM 主线如何一步步收敛到正式结果，不再把 baseline/current、性能口径和演示模式说明压在同一张流程图里。};

\begin{scope}[on background layer]
  \node[
    fit=(title)(subtitle)(n1)(s1)(n2)(s2)(n3)(s3)(n4)(s4)(n5)(s5)(footer),
    fill=panel,
    draw=linec,
    rounded corners=12pt,
    inner sep=16pt
  ] {};
\end{scope}

\end{tikzpicture}
\end{document}
